\begin{textsample}{POS Dim 3 – ap – Score 19.00 – ap/ap\_inf\_5.txt}  \label{ex:f3_pos_184}
II – \textbf{Crianças} bem \textbf{pequenas} ( 1 ano e 7 \textbf{meses} a 3 anos e 11 \textbf{meses} ) Nesta fase a \textbf{criança} irá aprender a utilizar adequadamente os seus sentidos, que a acompanham em todas as suas aventuras, na hora de explorar o mundo que a rodeia.
Isso irá permitir -lhe \textbf{assimilar} e perceber a informação que capta pelos seus sentidos, como diferenciar temperaturas, saber se faz frio ou calor.
Perceberá novas dimensões como o \textbf{afeto} ou o amor, pois as \textbf{crianças} descarregam emoções e refletem o mundo, assumindo diferentes posicionamentos ( SANTOS, 2001 ).
Iniciará a construção o seu pensamento por meio das suas ações e interações realizando imagens mentais com toda essa informação e vai avançando na expressão oral para \textbf{contar} as suas experiências.
Desde que aprende a falar irá fazê -lo a toda a hora, continuamente, mesmo que ninguém a \textbf{escute}.
A \textbf{criança} irá exigir a atenção dos outros enquanto fala, especialmente dos seus pais e ficará chateada se não for atendida.
As birras e chateações poderão ser uma constante nessa fase por conta de sua impaciência e egocentrismo.
Os pais e professores ao estipularem normas e limites lhes darão segurança e oportunidade para desenvolver valores, o bem e o mal criando -se o respeito, mas firme ao mesmo tempo nas situações.
É importante que o \textbf{adulto} esteja \textbf{atento} ao que surge durante os momentos de faz-de-conta para poder interferir de maneira a garantir a participação de todos e apontar caminhos que facilitem a integração do grupo ( SANTOS, 2001 ).
A \textbf{curiosidade} está aguçada, perguntando sempre o porquê das coisas.
Ela tem maior destreza física em geral participando com otimismo e segurança das atividades que envolvem o corpo.
Nestes \textbf{meses} veem -se tantas alterações na \textbf{criança} porque o cérebro cresce mais rapidamente durante os primeiros três anos de vida.
Cada \textbf{criança} é única e diferente, cada uma aprende ao seu ritmo, algumas mais rápido que outras.
Muitas têm problemas de alterações repentinas e há que dar -lhes tempo para se adaptarem às novas pessoas e lugares.
Quanto à interação, a \textbf{criança} amplia o seu mundo social, não se relacionando somente com os seus pais ou irmãos.
Mas, é necessário contatar com outras \textbf{crianças} : da escola, do \textbf{parque}, vizinhos, dentre outras, pois : > Simultaneamente, nesta etapa, as \textbf{crianças} tomam contato com o mundo que as cerca, através das experiências diretas com as pessoas e as coisas deste mundo e com as formas de expressão que nele ocorrem.
Esta inserção das \textbf{crianças} no mundo não seria possível sem que atividades voltadas simultaneamente para \textbf{cuidar} e educar estivessem presentes. ( BUJES, 2001, p. 16 ).
Com este interagir, a \textbf{criança} começa a tornar -se um ser social que compartilha e respeita as normas dos jogos.
O jogo irá proporcioná -la a capacidade de tomar as suas próprias decisões, como as de planificação e de construção onde acaba por desenvolver a sua criatividade, onde experimenta a sensação de domínio.
Uma boa ideia é incentivá -la quando quiser aprender alguma coisa, \textbf{emocionar} -se quando esta realizar ações por si só, etc.
Tudo isto fará com que ela se sinta bem e estimulará a continuação de aprender enquanto \textbf{brinca}. 1 ano e 7 \textbf{meses} - Imita \textbf{sons} e aumenta seu vocabulário, mas ainda não muito extenso ; - \textbf{Brinca} de espalhar e guardar tudo, claro que a seu modo ; - Normalmente já tem seu aparato motor pronto para andar, mas se ela não andar aos dez \textbf{meses} ou um ano e meio, ainda não é o caso de se preocupar, pois há diferenças individuais.
É importante que ela passe por todas as fases de se arrastar, segurar, trepar, engatinhar e andar. 2 anos - \textbf{Adultos} precisarão de bastante disponibilidade para responder a todos os questionamentos da \textbf{criança} : “ Como? ”, “ Quando? ” e o preferido “ Por quê? ” ; - Reafirmação da sua independência ; - Reconhece algumas cores e formas ; - Experiência do período sensório motor em que a aquisição do conhecimento acontece por meio dos sentidos. 3 anos - Coordenação fina está mais segura e é, geralmente, nessa época que a lateralidade ( destra ou canhota ) normalmente se define ; - Compreende perfeitamente o significado da palavra “ não ” e outras palavras de ordem ; - A linguagem oral permite -lhe falar com os outros com bastante seriedade ; - Descobre o prazer em \textbf{brincar} com o outro ; - \textbf{Assimila} centenas de palavras em pouco tempo ; - Reconhece e classifica formas, cores e espessuras. 3 anos e 11 \textbf{meses} - Exagera na cantoria, \textbf{brinca} com palavras, músicas e poesias ; - Nos \textbf{livros} identifica figuras ; - Atividades com cubo e lego ; - Ampliação o vocabulário ; - Habilidade manual cresce a cada \textbf{dia}, além de pegar \textbf{pequenos} objetos, folhear \textbf{livros} e abrir tampas de garrafas, ele já é capaz de utilizar a tesoura \textbf{infantil} ; - É cada vez mais independente ao nível da sua \textbf{higiene} ( desfralde ) e é já capaz de controlar os esfíncteres ( sobretudo durante o \textbf{dia} ).

% matched lemmas: adulto, afeto, assimilar, atento, brincar, contar, criança, cuidar, curiosidade, dia, emocionar, escutar, higiene, infantil, livro, mês, parque, pequeno, som
\end{textsample}
